% --- Archon physics formalization source begin ---
% archon:physics
% archon:covers PhyXMiniProblems/problem_phyx_mini_0097.lean
% archon:source-report reports/phyx_mini/problem_phyx_mini_0097.source.json

\chapter{Physics problem phyx\_mini\_0097}
\label{ch:PhyXMiniProblems_problem_phyx_mini_0097}

\paragraph{Problem source.}
\#\# Physical scenario

A ray of light traveling in air is incident at angle \$\textbackslash{}theta\_a\$ on one face of a 90.0° prism made of glass. Part of the light refracts into the prism and strikes the opposite face at point A (figure).The ray at A is at the critical angle.

\#\# Question

What is the value of \$\textbackslash{}theta\_a\$?

\#\# Answer choices

- A: A:  90°

- B: B:  60°

- C: C:  45°

- D: D:  135°

\#\# Figure caption (auxiliary; use the image as primary evidence)

The image depicts a right-angled triangle with a light ray entering and traversing through it. The triangle has a 90-degree angle marked at one vertex. The path of the light ray is shown with a purple line entering from one side of the triangle and exiting through point A on the opposite side.

There is a normal line (perpendicular to the side of the triangle) where the light ray enters. The angle between the incoming light ray and the normal is labeled \textbackslash{}(\textbackslash{}theta\_a\textbackslash{}). Additionally, the angle of incidence at the point of entry is marked as \textbackslash{}(40.0\textasciicircum{}\textbackslash{}circ\textbackslash{}). 

The right angle of the triangle is marked with \textbackslash{}(90.0\textasciicircum{}\textbackslash{}circ\textbackslash{}), indicating its perpendicular nature. The setup illustrates the concept of refraction and angle measurement in optics.

\#\# Dataset metadata

Geometrical Optics; Physical Model Grounding Reasoning, Multi-Formula Reasoning

\paragraph{Recorded answer/context.}
A: A:  90°

\paragraph{Figure/image path.}
/root/proposal\_for\_physic/science-mango/phyx\_mini\_run/phyx\_data/test\_image/97.png

\paragraph{Formalization target.}
create a compiling Lean file with sorry bodies at `PhyXMiniProblems/problem\_phyx\_mini\_0097.lean`. The Lean declarations must preserve the physical quantities, dimensions or dimensional roles, figure labels, governing-law hypotheses, and final relation expressed by this problem.
Use Mathlib/Physlib names found through LeanExplore where available. If a physics API is missing, introduce faithful local abstractions rather than scalar placeholder aliases.

\begin{theorem}[Physics formalization target]
\label{thm:physics:phyx_mini_0097:target}
\lean{PhyXMiniProblems.ProblemPhyXMini0097.problem_phyx_mini_0097}
\uses{def:physics:phyx-mini-0097:phyxminiproblems-problemphyxmini0097-degrees, def:physics:phyx-mini-0097:phyxminiproblems-problemphyxmini0097-rightangleprismraysetup, def:physics:phyx-mini-0097:phyxminiproblems-problemphyxmini0097-thetaa, def:physics:phyx-mini-0097:phyxminiproblems-problemphyxmini0097-hasphysicalparameters, def:physics:phyx-mini-0097:phyxminiproblems-problemphyxmini0097-matchesproblemandfigure, def:physics:phyx-mini-0097:phyxminiproblems-problemphyxmini0097-satisfiesrightangleprismgeometry, def:physics:phyx-mini-0097:phyxminiproblems-problemphyxmini0097-obeyssnellrefraction, def:physics:phyx-mini-0097:phyxminiproblems-problemphyxmini0097-isatcriticalangleata}
The assigned autoformalize agent should translate this physics problem into Lean declarations in the covered file, with theorem and lemma proof bodies written as `by sorry`.
\end{theorem}
\begin{proof}
This is an autoformalization task, not a proof task. Produce statements that can later be proved by the physics prover without weakening the physical model.
\end{proof}

% --- Archon declaration topology begin ---
\section{Declaration topology}
\begin{definition}[Optical Medium]
  \label{def:physics:phyx-mini-0097:phyxminiproblems-problemphyxmini0097-opticalmedium}
  \lean{PhyXMiniProblems.ProblemPhyXMini0097.OpticalMedium}
  The two homogeneous optical media traversed by the pictured ray
\end{definition}
\begin{proof}
  This is the declared model definition of Optical Medium.
\end{proof}
\begin{definition}[Prism Face]
  \label{def:physics:phyx-mini-0097:phyxminiproblems-problemphyxmini0097-prismface}
  \lean{PhyXMiniProblems.ProblemPhyXMini0097.PrismFace}
  The three faces in the triangular prism cross-section
\end{definition}
\begin{proof}
  This is the declared model definition of Prism Face.
\end{proof}
\begin{definition}[Prism Vertex]
  \label{def:physics:phyx-mini-0097:phyxminiproblems-problemphyxmini0097-prismvertex}
  \lean{PhyXMiniProblems.ProblemPhyXMini0097.PrismVertex}
  The three vertices in the triangular prism cross-section
\end{definition}
\begin{proof}
  This is the declared model definition of Prism Vertex.
\end{proof}
\begin{definition}[Figure Point]
  \label{def:physics:phyx-mini-0097:phyxminiproblems-problemphyxmini0097-figurepoint}
  \lean{PhyXMiniProblems.ProblemPhyXMini0097.FigurePoint}
  The two ray-interaction points distinguished in the figure
\end{definition}
\begin{proof}
  This is the declared model definition of Figure Point.
\end{proof}
\begin{definition}[face]
  \label{def:physics:phyx-mini-0097:phyxminiproblems-problemphyxmini0097-figurepoint-face}
  \lean{PhyXMiniProblems.ProblemPhyXMini0097.FigurePoint.face}
  \uses{def:physics:phyx-mini-0097:phyxminiproblems-problemphyxmini0097-prismface, def:physics:phyx-mini-0097:phyxminiproblems-problemphyxmini0097-figurepoint}
  The prism face on which each distinguished interaction point lies
\end{definition}
\begin{proof}
  This is the declared model definition of face.
\end{proof}
\begin{definition}[Ray Segment]
  \label{def:physics:phyx-mini-0097:phyxminiproblems-problemphyxmini0097-raysegment}
  \lean{PhyXMiniProblems.ProblemPhyXMini0097.RaySegment}
  The three consecutive directed portions of the optical path
\end{definition}
\begin{proof}
  This is the declared model definition of Ray Segment.
\end{proof}
\begin{definition}[medium]
  \label{def:physics:phyx-mini-0097:phyxminiproblems-problemphyxmini0097-raysegment-medium}
  \lean{PhyXMiniProblems.ProblemPhyXMini0097.RaySegment.medium}
  \uses{def:physics:phyx-mini-0097:phyxminiproblems-problemphyxmini0097-opticalmedium, def:physics:phyx-mini-0097:phyxminiproblems-problemphyxmini0097-raysegment}
  The physical medium occupied by each portion of the optical path
\end{definition}
\begin{proof}
  This is the declared model definition of medium.
\end{proof}
\begin{definition}[Prism Interface]
  \label{def:physics:phyx-mini-0097:phyxminiproblems-problemphyxmini0097-prisminterface}
  \lean{PhyXMiniProblems.ProblemPhyXMini0097.PrismInterface}
  The two transparent interfaces encountered by the ray, in path order
\end{definition}
\begin{proof}
  This is the declared model definition of Prism Interface.
\end{proof}
\begin{definition}[point]
  \label{def:physics:phyx-mini-0097:phyxminiproblems-problemphyxmini0097-prisminterface-point}
  \lean{PhyXMiniProblems.ProblemPhyXMini0097.PrismInterface.point}
  \uses{def:physics:phyx-mini-0097:phyxminiproblems-problemphyxmini0097-figurepoint, def:physics:phyx-mini-0097:phyxminiproblems-problemphyxmini0097-prisminterface}
  The labelled point at which an interface crossing occurs
\end{definition}
\begin{proof}
  This is the declared model definition of point.
\end{proof}
\begin{definition}[incident Segment]
  \label{def:physics:phyx-mini-0097:phyxminiproblems-problemphyxmini0097-prisminterface-incidentsegment}
  \lean{PhyXMiniProblems.ProblemPhyXMini0097.PrismInterface.incidentSegment}
  \uses{def:physics:phyx-mini-0097:phyxminiproblems-problemphyxmini0097-raysegment, def:physics:phyx-mini-0097:phyxminiproblems-problemphyxmini0097-prisminterface}
  The segment incident on a given interface
\end{definition}
\begin{proof}
  This is the declared model definition of incident Segment.
\end{proof}
\begin{definition}[transmitted Segment]
  \label{def:physics:phyx-mini-0097:phyxminiproblems-problemphyxmini0097-prisminterface-transmittedsegment}
  \lean{PhyXMiniProblems.ProblemPhyXMini0097.PrismInterface.transmittedSegment}
  \uses{def:physics:phyx-mini-0097:phyxminiproblems-problemphyxmini0097-raysegment, def:physics:phyx-mini-0097:phyxminiproblems-problemphyxmini0097-prisminterface}
  The limiting or actual transmitted segment at a given interface
\end{definition}
\begin{proof}
  This is the declared model definition of transmitted Segment.
\end{proof}
\begin{definition}[Boundary Behavior]
  \label{def:physics:phyx-mini-0097:phyxminiproblems-problemphyxmini0097-boundarybehavior}
  \lean{PhyXMiniProblems.ProblemPhyXMini0097.BoundaryBehavior}
  The qualitative refraction regime shown at a transparent interface
\end{definition}
\begin{proof}
  This is the declared model definition of Boundary Behavior.
\end{proof}
\begin{definition}[degrees]
  \label{def:physics:phyx-mini-0097:phyxminiproblems-problemphyxmini0097-degrees}
  \lean{PhyXMiniProblems.ProblemPhyXMini0097.degrees}
  Convert a numerical degree readout to Mathlib's physical angle type
\end{definition}
\begin{proof}
  This is the declared model definition of degrees.
\end{proof}
\begin{definition}[Right Angle Prism Ray Setup]
  \label{def:physics:phyx-mini-0097:phyxminiproblems-problemphyxmini0097-rightangleprismraysetup}
  \lean{PhyXMiniProblems.ProblemPhyXMini0097.RightAnglePrismRaySetup}
  \uses{def:physics:phyx-mini-0097:phyxminiproblems-problemphyxmini0097-opticalmedium, def:physics:phyx-mini-0097:phyxminiproblems-problemphyxmini0097-prismface, def:physics:phyx-mini-0097:phyxminiproblems-problemphyxmini0097-prismvertex, def:physics:phyx-mini-0097:phyxminiproblems-problemphyxmini0097-raysegment, def:physics:phyx-mini-0097:phyxminiproblems-problemphyxmini0097-prisminterface, def:physics:phyx-mini-0097:phyxminiproblems-problemphyxmini0097-boundarybehavior}
  The material data and angular readouts for the pictured prism and ray. angleToNormal segment face records a genuine normal-relative angle, rather than replacing the light ray or prism by a scalar. Only the combinations selected below by the two interfaces have a role in the physical laws
\end{definition}
\begin{proof}
  This is the declared model definition of Right Angle Prism Ray Setup.
\end{proof}
\begin{definition}[theta A]
  \label{def:physics:phyx-mini-0097:phyxminiproblems-problemphyxmini0097-thetaa}
  \lean{PhyXMiniProblems.ProblemPhyXMini0097.thetaA}
  \uses{def:physics:phyx-mini-0097:phyxminiproblems-problemphyxmini0097-rightangleprismraysetup}
  The angle labelled θₐ, measured in air from the entry-face normal
\end{definition}
\begin{proof}
  This is the declared model definition of theta A.
\end{proof}
\begin{definition}[entry Refraction Angle]
  \label{def:physics:phyx-mini-0097:phyxminiproblems-problemphyxmini0097-entryrefractionangle}
  \lean{PhyXMiniProblems.ProblemPhyXMini0097.entryRefractionAngle}
  \uses{def:physics:phyx-mini-0097:phyxminiproblems-problemphyxmini0097-rightangleprismraysetup}
  The refracted angle in the glass immediately after the entry face
\end{definition}
\begin{proof}
  This is the declared model definition of entry Refraction Angle.
\end{proof}
\begin{definition}[incidence Angle At A]
  \label{def:physics:phyx-mini-0097:phyxminiproblems-problemphyxmini0097-incidenceangleata}
  \lean{PhyXMiniProblems.ProblemPhyXMini0097.incidenceAngleAtA}
  \uses{def:physics:phyx-mini-0097:phyxminiproblems-problemphyxmini0097-rightangleprismraysetup}
  The internal incidence angle at point A, measured from its face normal
\end{definition}
\begin{proof}
  This is the declared model definition of incidence Angle At A.
\end{proof}
\begin{definition}[transmission Angle At A]
  \label{def:physics:phyx-mini-0097:phyxminiproblems-problemphyxmini0097-transmissionangleata}
  \lean{PhyXMiniProblems.ProblemPhyXMini0097.transmissionAngleAtA}
  \uses{def:physics:phyx-mini-0097:phyxminiproblems-problemphyxmini0097-rightangleprismraysetup}
  The limiting transmitted angle in air at A, measured from the normal
\end{definition}
\begin{proof}
  This is the declared model definition of transmission Angle At A.
\end{proof}
\begin{definition}[Is Physical Normal Angle]
  \label{def:physics:phyx-mini-0097:phyxminiproblems-problemphyxmini0097-isphysicalnormalangle}
  \lean{PhyXMiniProblems.ProblemPhyXMini0097.IsPhysicalNormalAngle}
  A normal-relative ray angle on the physical branch from 0° to 90°
\end{definition}
\begin{proof}
  This is the declared model definition of Is Physical Normal Angle.
\end{proof}
\begin{definition}[Has Physical Parameters]
  \label{def:physics:phyx-mini-0097:phyxminiproblems-problemphyxmini0097-hasphysicalparameters}
  \lean{PhyXMiniProblems.ProblemPhyXMini0097.HasPhysicalParameters}
  \uses{def:physics:phyx-mini-0097:phyxminiproblems-problemphyxmini0097-prisminterface, def:physics:phyx-mini-0097:phyxminiproblems-problemphyxmini0097-rightangleprismraysetup, def:physics:phyx-mini-0097:phyxminiproblems-problemphyxmini0097-isphysicalnormalangle}
  Positivity, optical ordering, and physical angle branches for the setup
\end{definition}
\begin{proof}
  This is the declared model definition of Has Physical Parameters.
\end{proof}
\begin{definition}[Matches Problem And Figure]
  \label{def:physics:phyx-mini-0097:phyxminiproblems-problemphyxmini0097-matchesproblemandfigure}
  \lean{PhyXMiniProblems.ProblemPhyXMini0097.MatchesProblemAndFigure}
  \uses{def:physics:phyx-mini-0097:phyxminiproblems-problemphyxmini0097-degrees, def:physics:phyx-mini-0097:phyxminiproblems-problemphyxmini0097-rightangleprismraysetup, def:physics:phyx-mini-0097:phyxminiproblems-problemphyxmini0097-entryrefractionangle}
  Literal data from the problem and image: the apex between the two crossed faces is a right angle, the entry normal is shown, the first crossing is an ordinary partial refraction, and the in-glass entry angle is 40.0°
\end{definition}
\begin{proof}
  This is the declared model definition of Matches Problem And Figure.
\end{proof}
\begin{definition}[Satisfies Right Angle Prism Geometry]
  \label{def:physics:phyx-mini-0097:phyxminiproblems-problemphyxmini0097-satisfiesrightangleprismgeometry}
  \lean{PhyXMiniProblems.ProblemPhyXMini0097.SatisfiesRightAnglePrismGeometry}
  \uses{def:physics:phyx-mini-0097:phyxminiproblems-problemphyxmini0097-rightangleprismraysetup, def:physics:phyx-mini-0097:phyxminiproblems-problemphyxmini0097-entryrefractionangle, def:physics:phyx-mini-0097:phyxminiproblems-problemphyxmini0097-incidenceangleata}
  Plane geometry for a ray crossing two prism faces meeting at the apex: its two internal normal-relative angles add to the angle between those faces
\end{definition}
\begin{proof}
  This is the declared model definition of Satisfies Right Angle Prism Geometry.
\end{proof}
\begin{definition}[Satisfies Snell Law At]
  \label{def:physics:phyx-mini-0097:phyxminiproblems-problemphyxmini0097-satisfiessnelllawat}
  \lean{PhyXMiniProblems.ProblemPhyXMini0097.SatisfiesSnellLawAt}
  \uses{def:physics:phyx-mini-0097:phyxminiproblems-problemphyxmini0097-prisminterface, def:physics:phyx-mini-0097:phyxminiproblems-problemphyxmini0097-rightangleprismraysetup}
  Snell's law at either transparent prism face: refractive index times the sine of the normal-relative ray angle is conserved across the interface
\end{definition}
\begin{proof}
  This is the declared model definition of Satisfies Snell Law At.
\end{proof}
\begin{definition}[Obeys Snell Refraction]
  \label{def:physics:phyx-mini-0097:phyxminiproblems-problemphyxmini0097-obeyssnellrefraction}
  \lean{PhyXMiniProblems.ProblemPhyXMini0097.ObeysSnellRefraction}
  \uses{def:physics:phyx-mini-0097:phyxminiproblems-problemphyxmini0097-rightangleprismraysetup, def:physics:phyx-mini-0097:phyxminiproblems-problemphyxmini0097-satisfiessnelllawat}
  Snell refraction governs both transparent boundary interactions
\end{definition}
\begin{proof}
  This is the declared model definition of Obeys Snell Refraction.
\end{proof}
\begin{definition}[Is At Critical Angle At A]
  \label{def:physics:phyx-mini-0097:phyxminiproblems-problemphyxmini0097-isatcriticalangleata}
  \lean{PhyXMiniProblems.ProblemPhyXMini0097.IsAtCriticalAngleAtA}
  \uses{def:physics:phyx-mini-0097:phyxminiproblems-problemphyxmini0097-degrees, def:physics:phyx-mini-0097:phyxminiproblems-problemphyxmini0097-rightangleprismraysetup, def:physics:phyx-mini-0097:phyxminiproblems-problemphyxmini0097-transmissionangleata}
  At A the ray is at the critical incidence: the limiting transmitted ray is tangent to the opposite face and hence is 90° from that face's normal
\end{definition}
\begin{proof}
  This is the declared model definition of Is At Critical Angle At A.
\end{proof}
\begin{definition}[Answer Choice]
  \label{def:physics:phyx-mini-0097:phyxminiproblems-problemphyxmini0097-answerchoice}
  \lean{PhyXMiniProblems.ProblemPhyXMini0097.AnswerChoice}
  Labels of the four multiple-choice answers printed with the problem
\end{definition}
\begin{proof}
  This is the declared model definition of Answer Choice.
\end{proof}
\begin{definition}[answer Angle Degrees]
  \label{def:physics:phyx-mini-0097:phyxminiproblems-problemphyxmini0097-answerangledegrees}
  \lean{PhyXMiniProblems.ProblemPhyXMini0097.answerAngleDegrees}
  \uses{def:physics:phyx-mini-0097:phyxminiproblems-problemphyxmini0097-answerchoice}
  Numerical degree readout printed beside each answer label
\end{definition}
\begin{proof}
  This is the declared model definition of answer Angle Degrees.
\end{proof}
\begin{definition}[recorded Answer Choice]
  \label{def:physics:phyx-mini-0097:phyxminiproblems-problemphyxmini0097-recordedanswerchoice}
  \lean{PhyXMiniProblems.ProblemPhyXMini0097.recordedAnswerChoice}
  \uses{def:physics:phyx-mini-0097:phyxminiproblems-problemphyxmini0097-answerchoice}
  The answer label stored in the dataset, retained only as source metadata
\end{definition}
\begin{proof}
  This is the declared model definition of recorded Answer Choice.
\end{proof}
% --- Archon declaration topology end ---
% --- Archon physics formalization source end ---
